(1) 记两个单界面 Fresnel 反射幅度
\[
r_{01}=\frac{n_0-n_1}{n_0+n_1},\qquad r_{12}=\frac{n_1-n_s}{n_1+n_s},
\]
膜单程相位 $\delta=2\pi n_1d/\lambda$。利用多次往返传播 $e^{2i\delta}$ 与界面振幅关系，整体复反射幅度为
\[
r(\lambda)=\frac{r_{01}+r_{12}e^{2i\delta}}{1+r_{01}r_{12}e^{2i\delta}},
\qquad \mathcal R=|r|^2.
\]
这是电场幅度的相干和，不可以先将各反射强度相加。

(2) 四分之一波条件 $\delta(\lambda_0)=\pi/2$ 使 $e^{2i\delta}=-1$。此时 $r=(r_{01}-r_{12})/(1-r_{01}r_{12})$，所以零反射要求 $r_{01}=r_{12}$。化简得
\[
n_1^2=n_0n_s,\qquad n_1=\sqrt{n_0n_s}.
\]
使 $\delta_0=\pi/2$ 的最小正厚度为
\[
d_{\min}=\frac{\lambda_0}{4\sqrt{n_0n_s}}.
\]
因为 $n_s>n_0$，此时 $n_0<n_1<n_s$，界面反射幅度均为负数，恰以半波往返相位抵消。

(3) 在零反射设计下令 $\rho=r_{01}=r_{12}$，$|\rho|<1$，以及 $\varepsilon=\delta-\pi/2$。有 $e^{2i\delta}=-e^{2i\varepsilon}$，因此
\[
r=\frac{\rho(1-e^{2i\varepsilon})}{1-\rho^2e^{2i\varepsilon}}
=-\frac{2i\rho}{1-\rho^2}\varepsilon+O(\varepsilon^2).
\]
由 $\delta=(\pi/2)(\lambda_0/\lambda)$，$\varepsilon=-(\pi/2)(\Delta\lambda/\lambda_0)+O((\Delta\lambda/\lambda_0)^2)$，所以
\[
\mathcal R(\lambda)
=\frac{\pi^2\rho^2}{(1-\rho^2)^2}
\left(\frac{\Delta\lambda}{\lambda_0}\right)^2
+O\!\left(\left|\frac{\Delta\lambda}{\lambda_0}\right|^3\right).
\]
一阶反射率项消失，说明设计中心处是严格的反射率极小值。

(4) 正入射时能流与介质折射率成正比，若整体透射电场幅度为 $t$，则 $\mathcal T=(n_s/n_0)|t|^2$，不可直接用 $|t|^2$ 代替能流透射率。无吸收多层膜满足 $\mathcal R+\mathcal T=1$，所以设计波长 $\mathcal R=0,\mathcal T=1$。偏离设计波长、入射角非零、膜吸收或有色散时不再保证完全增透。
